\hypertarget{the-pyobject-core-object-model}{%
\chapter{THE PYOBJECT CORE OBJECT
MODEL}\label{the-pyobject-core-object-model}}

\hypertarget{the-unified-pyobject-pyvarobject-structs}{%
\subsection{2.1 The Unified PyObject \& PyVarObject
Structs}\label{the-unified-pyobject-pyvarobject-structs}}

In CPython, \textbf{every python value is a heap-allocated C struct}.
There are no primitive variables on the stack. The fundamental base
structure is defined in \passthrough{\lstinline!Include/object.h!}:

\begin{lstlisting}[language=C]
typedef struct _object {
    _PyObject_HEAD_EXTRA    // Doubly-linked list pointers for active heap tracing
    Py_ssize_t ob_refcnt;   // Reference count counter
    struct _typeobject *ob_type; // Pointer to the type definition object
} PyObject;

typedef struct {
    PyObject ob_base;       // Base type fields
    Py_ssize_t ob_size;     // Size of the variable portion (e.g. string length)
} PyVarObject;
\end{lstlisting}

\hypertarget{comparison-to-c-memory}{%
\subsubsection{Comparison to C++ memory:}\label{comparison-to-c-memory}}

Unlike C++, where objects are laid out contiguously on the stack or heap
based on their static type declarations and resolved via virtual
function tables (\passthrough{\lstinline!vtable!}), CPython objects are
entirely dynamic: 1. \textbf{Uniform Pointer Width}: Every variable in
Python is a pointer (\passthrough{\lstinline!PyObject*!}), consuming
exactly 8 bytes (on 64-bit platforms). 2. \textbf{No VTable
Indirection}: Polymorphism is resolved dynamically by dereferencing the
type pointer \passthrough{\lstinline!ob\_type!} at runtime to lookup
operations.

\begin{lstlisting}
Python Reference Pointer:
[Name: x] ---> [ Heap Allocation: PyObject ]
               +--------------------------------------+
               | _PyObject_HEAD_EXTRA (Tracking)      |
               | ob_refcnt: 1                         |
               | ob_type: ----> [ PyTypeObject (int)] |
               | Integer Data Value Payload           |
               +--------------------------------------+
\end{lstlisting}

\hypertarget{the-type-object-slot-system}{%
\subsection{2.2 The Type-Object Slot
System}\label{the-type-object-slot-system}}

The type of an object is defined by
\passthrough{\lstinline!PyTypeObject!} (which is itself a subclass of
\passthrough{\lstinline!PyObject!}). * \textbf{Type Slots}: The type
struct defines a series of ``slots'' (pointers to C functions) that
determine how the object behaves when subjected to operators. -
\passthrough{\lstinline!tp\_alloc!}: Allocates heap memory for the
struct. - \passthrough{\lstinline!tp\_new!}: Constructor initializer
(invokes allocation). - \passthrough{\lstinline!tp\_init!}:
Initialization hook (corresponds to
\passthrough{\lstinline!\_\_init\_\_!}). -
\passthrough{\lstinline!tp\_dealloc!}: Destructor wrapper, called when
reference counts hit zero. - \passthrough{\lstinline!tp\_hash!}: Defines
the hashing method for dictionaries and sets.

\begin{lstlisting}[language=Python]
# Customizing slots via class definitions
class CustomObject:
    def __init__(self, val):
        self.val = val

    def __hash__(self):
        # Maps to the tp_hash C slot
        return hash(self.val)

obj = CustomObject("test")
print("Hash via type slot:", hash(obj))
\end{lstlisting}

\hypertarget{reference-counting-lifecycle}{%
\subsection{2.3 Reference Counting
Lifecycle}\label{reference-counting-lifecycle}}

CPython manages memory directly via reference counting. The macros
\passthrough{\lstinline!Py\_INCREF()!} and
\passthrough{\lstinline!Py\_DECREF()!} increment and decrement
\passthrough{\lstinline!ob\_refcnt!} respectively. * \textbf{The
Decrement Process}: When \passthrough{\lstinline!Py\_DECREF(obj)!} drops
the reference count to zero: 1. The runtime invokes
\passthrough{\lstinline!obj->ob\_type->tp\_dealloc(obj)!}. 2. This
deallocator function decrements references of any nested objects inside
this object. 3. The raw memory is released back to the interpreter's
memory allocator (PyMalloc). * \textbf{Leak Debugging}: In custom C
extensions, failing to decrement references leads to memory leaks. In
debug builds, compiling with \passthrough{\lstinline!Py\_TRACE\_REFS!}
tracks all active references in a global list.

\begin{lstlisting}[language=Python]
import sys

# Reference counts tracking
my_list = [100, 200]
print("Starting references count:", sys.getrefcount(my_list) - 1)

ref_a = my_list
ref_b = my_list
print("References after bindings:", sys.getrefcount(my_list) - 1)

del ref_a
print("References after deleting one binding:", sys.getrefcount(my_list) - 1)
\end{lstlisting}

\hypertarget{built-in-caching-and-interning-optimizations}{%
\subsection{2.4 Built-in Caching and Interning
Optimizations}\label{built-in-caching-and-interning-optimizations}}

To avoid allocating memory for common objects, CPython uses global
caches: * \textbf{Small Integer Cache}: CPython pre-allocates an array
of integer objects from \passthrough{\lstinline!-5!} to
\passthrough{\lstinline!256!} (\passthrough{\lstinline!NSMALLNEGINTS!}
and \passthrough{\lstinline!NSMALLPOSINTS!}). - \emph{Under the hood}:
Any assignment in this range returns a pointer to the pre-existing
static object, bypassing heap allocations. * \textbf{String Interning}:
String literals resembling identifiers are automatically
\textbf{interned}. - \emph{Mechanism}: Interned strings are stored in a
global dictionary. If a string is already interned, any new
instantiation returns the existing string pointer. This allows matching
strings via fast pointer comparison (\(O(1)\) address checks) instead of
byte-by-byte comparison (\(O(N)\) string scans).

\begin{lstlisting}[language=Python]
# Small Integer Caching Check
int_a = 256
int_b = 256
print("Are 256 pointers identical?", int_a is int_b) # True (cached)

int_c = 257
int_d = 257
print("Are 257 pointers identical?", int_c is int_d) # False (dynamically allocated)

# String Interning Check
str_a = "godhood_string"
str_b = "godhood_string"
print("Are literal strings interned?", str_a is str_b) # True
\end{lstlisting}

\begin{center}\rule{0.5\linewidth}{0.5pt}\end{center}
